%% SYNTHETIC TEST UNIT generated by writing/tools/make_outline_book.py from FINAL_OUTLINE_v2 -- not book text
\chapter{幂次修正、次领头幂次、重整化子与非微扰地位}\label{ch:power}
本单元为测试用的合成文本（由 FINAL\_OUTLINE\_v2 自动生成），陈述内容从略。

\section{第 1 节（测试）}\label{sec:power-twist4}

\begin{proposition}[DIS 扭度 4 修正的 $x$ 空间形式 \tagA]\label{prop:power-dis-twist4-x}
（测试用合成陈述：内容从略。）
非微扰地位：\cref{asm:intro-bridge}。
\uses{prop:dis-twist4, thm:disall-factorization, asm:intro-bridge}
\end{proposition}
\begin{proofsketch}测试。\end{proofsketch}
\paragraph{地位与缺口}测试：本陈述的地位由\cref{conj:power-renormalon-ht}、\cref{prop:power-rearrangement}决定。

\begin{proposition}[SIDIS 型过程中的扭度 3 共线因子化 \tagP]\label{prop:power-twist3-sidis}
（测试用合成陈述：内容从略。）
\uses{def:tmd-etqs, prop:tmd-twist3-matching}
\end{proposition}
\paragraph{地位与缺口}测试：见文献。

\begin{conjecture}[强子—强子碰撞中单举强子单自旋不对称的扭度 3 因子化 \tagC]\label{conj:power-twist3-pp}
（测试用合成陈述：内容从略。）
\uses{prop:power-twist3-sidis, conj:dy-coloured-final}
\end{conjecture}
\paragraph{地位与缺口}测试：本陈述的地位由\cref{conj:dy-coloured-final}决定。

\section{第 2 节（测试）}\label{sec:power-twist3}

\begin{example}[最低阶的软胶子极点机制]\label{ex:power-sgp}
（测试用合成陈述：内容从略。）
\uses{prop:power-twist3-sidis, thm:tmd-sivers-sign}
\end{example}

\begin{proposition}[次软定理 \tagR]\label{prop:power-next-to-soft}
（$R_{\mathrm{PT}}$：在重整化微扰论中对红外正规化的部分子态严格。）
（测试用合成陈述：内容从略。）
\uses{prop:qcd-soft-collinear-limits}
\end{proposition}

\begin{proposition}[次领头幂次因子化与端点发散 \tagP]\label{prop:power-nlp-endpoint}
（测试用合成陈述：内容从略。）
\uses{prop:power-next-to-soft, prop:thr-nlp}
\end{proposition}
\paragraph{地位与缺口}测试：本陈述的地位由\cref{prop:heavy-qcdf-endpoint}决定。

\section{第 3 节（测试）}\label{sec:power-nlp}

\begin{proposition}[OPE 基的重排自由 \tagR]\label{prop:power-rearrangement}
（$R_{\mathrm{def}}$：由算符定义直接推出。）
（测试用合成陈述：内容从略。）
\uses{def:dis-twist}
\end{proposition}

\begin{remark}["实用 OPE" 与重整化子吸收的两个猜想性成分]\label{rem:power-practical-ope}
（测试用合成陈述：内容从略。）
% uses: none
\end{remark}

\begin{definition}[Borel 变换与重整化子]\label{def:power-borel}
（测试用合成陈述：内容从略。）
\uses{prop:prelim-beta, def:qcd-lambda}
\end{definition}

\section{第 4 节（测试）}\label{sec:power-ope-basis}

\begin{proposition}[大 $\beta_0$ 极限下的 Adler 函数 \tagP]\label{prop:power-adler-large-beta0}
（测试用合成陈述：内容从略。）
\uses{def:power-borel, def:ee-adler, thm:ee-ope}
\end{proposition}
\paragraph{地位与缺口}测试：见文献。

\begin{example}[带 $n$ 个费米子泡的单胶子]\label{ex:power-renormalon-chain}
（测试用合成陈述：内容从略。）
\uses{prop:power-adler-large-beta0}
\end{example}

\begin{proposition}[若干过程中的重整化子与幂次修正 \tagP]\label{prop:power-renormalon-examples}
（测试用合成陈述：内容从略。）
\uses{def:power-borel, prop:power-adler-large-beta0}
\end{proposition}
\paragraph{地位与缺口}测试：见文献。

\section{第 5 节（测试）}\label{sec:power-renormalons}

\begin{conjecture}[红外重整化子与高扭度的一般对应 \tagC]\label{conj:power-renormalon-ht}
（测试用合成陈述：内容从略。）
\uses{prop:power-renormalon-examples, rem:power-practical-ope}
\end{conjecture}
\paragraph{地位与缺口}测试：本陈述的地位由\cref{asm:intro-np-ope}决定。

\begin{remark}[非微扰地位]\label{rem:power-np-status}
（测试用合成陈述：内容从略。）
非微扰地位：\cref{rem:power-np-status}。
% uses: none
\end{remark}

\begin{theorem}[微扰 $\varphi^4_4$ 中 OPE 的收敛性 \tagR]\label{thm:power-ope-convergence}
（$R_{\mathrm{PT}}$：在重整化微扰论中对红外正规化的部分子态严格。）
（测试用合成陈述：内容从略。）
% uses: none
\end{theorem}

\section{第 6 节（测试）}\label{sec:power-np}

\begin{proposition}['t Hooft 模型 \tagP]\label{prop:power-thooft-model}
（测试用合成陈述：内容从略。）
% uses: none
\end{proposition}
\paragraph{地位与缺口}测试：见文献。
